\begin{claim}{AC.27}{Proved; exact Cartier--Mumford rational test}
Let $k=\mathbf F_q$, $q=2^d$, and let
\[
 C/k:\quad Y^2+h(x)Y=f(x),\qquad \deg h\le2,
 \quad f\text{ monic of degree }5
\]
be smooth of genus two, based at its unique point $P=\infty$.
Let $n>1$ be odd and let $a\in J_C(\overline k)$ have exact order
$n$ and satisfy $\sigma_q(a)=[q]a$, where $\sigma_q$ is arithmetic
Frobenius. Choose a finite extension $E/k$ containing $a$ and $\mu_n$.
For $0<i<n$, write the unique reduced Mumford representative of $ia$
as $(u_i,v_i)$, with $u_i$ monic of degree $r_i\in\{1,2\}$, so that
$ia=[D_i-r_iP]$.

For $j=2i\bmod n$, compute $c_i$ as follows.
If $r_i=1$, put $c_i=h_1$, the coefficient of $x$ in $h$.
If $r_i=2$ and $\gcd(u_i,h)$ has degree one, put $c_i=h_2$.
Otherwise $r_i=2$ and $\gcd(u_i,h)=1$; let $v$ be the Cantor
doubling interpolation of degree at most three, with
$u_i^2\mid v^2+hv+f$ and $v\equiv v_i\bmod u_i$.
If the reduced double has degree one, put $c_i=1$.
If it has degree two, let $v_3$ be the cubic coefficient of $v$
and $b_j$ the coefficient of $x$ in $u_j$, and put $c_i=v_3b_j$.
No other case occurs for a nonzero point of odd order. Define
\[
 A_i=\prod_{s=0}^{d-1}c_{2^s i}^{\,2^{d-1-s}}\in E.       \tag{AC.14}
\]
There is a unique constant cyclic $P$-split degree-$n$ cover over $k$
with geometric character subgroup generated by $a$.
Its linear coherent $q$-Frobenius scalar on the summand for $ia$ is
$\sigma_q^{-1}(A_i)$. Consequently its relative class number is odd
if and only if
\[
                         A_i\ne1\qquad(0<i<n).           \tag{AC.15}
\]
For the primitive factor alone use only $\gcd(i,n)=1$.
These formulas refer to the pointing at infinity; another base point
requires the twist of AC.6.
\end{claim}
\begin{proof}
Put $L_i=\mathcal O(D_i-r_iP)$ and rigidify it at $P$ using
$t=x^2/Y$. Monicity of $f$ gives leading residues one for $xt^2$
and $Yt^5$. The differential $dx/h$ has divisor $2P$.
At finite unramified points this follows from $x$ being a local
parameter. At a finite ramification point,
$dx/h=dY/(f'+h'Y)$ is nonvanishing by smoothness.
Thus its canonical divisor is supported at $P$ and has degree two.
A basis of the one-dimensional Serre dual space
$H^0(C_E,\Omega_C\otimes L_i^{-1})$ is
\[
                            \omega_i=u_i\,dx/h.
\]
Indeed $\operatorname{div}(u_i)=D_i+\iota D_i-2r_iP$ and $r_i\le2$.
Let $e_i$ be its dual basis in $H^1(C_E,L_i)$.

Let $g_i$ have divisor $2D_i-D_j-(2r_i-r_j)P$, normalized to
leading residue one at $P$ relative to $t$.
Multiplication by $g_i$ identifies $L_i^2$ with $L_j$ respecting
rigidifications. Absolute Frobenius followed by this identification
is square-semilinear. Cartier is its Serre adjoint:
\[
 \langle g_iF(z),\omega\rangle
             =\langle z,\mathcal C(g_i\omega)\rangle^2.
\]
This identity follows directly in the residue pairing: the coefficient
of $t^{-1}dt$ in $z^2\eta$ is the square of the residue of
$z\mathcal C(\eta)$. Therefore
$\mathcal C(g_i\omega_j)=\sqrt{c_i}\omega_i$ is exactly the assertion
$F(e_i)=c_i e_j$.

We verify all four cases.
If $r_i=1$, the point is unramified, since a ramification point would
give a two-torsion class. Then $D_j=2D_i$, $u_j=u_i^2$, and $g_i=1$.
The identity $\mathcal C(dx/h)=\sqrt{h_1}\,dx/h$ gives $c_i=h_1$.

If $r_i=2$ and $\gcd(u_i,h)=x-w$, write $D_i=W+R$ with $W$
ramified and $R$ unramified. Since $2W\sim2P$, we have $D_j=2R$
and $g_i=x-w$. With $u_i=(x-w)(x-r)$, direct extraction of odd
powers gives
\[
 \mathcal C\bigl((x-w)(x-r)^2dx/h\bigr)
             =\sqrt{h_2}(x-w)(x-r)dx/h.
\]
For this calculation write $(x-w)/h=(x-w)h/h^2$ and use $h(w)=0$:
the square root of the odd part is $\sqrt{h_2}(x-w)$.
Hence $c_i=h_2$. Two ramification points in $D_i$ would give a
two-torsion class and are excluded.

In the remaining cases $\gcd(u_i,h)=1$.
For a degree-two reduced double, Cantor reduction gives
\[
 v^2+hv+f=v_3^2u_i^2u_j,\qquad
                     g_i=(Y+v)/(v_3u_j).
\]
This function has the required divisor and normalized leading residue.
The Artin--Schreier relation for $Y/h$ gives
\[
 \mathcal C((Y+v)dx/h)=v_3\sqrt{b_j}\,u_i dx/h.
\]
Indeed the term $(Y/h)^2dx$ has Cartier zero, and the odd powers of
$f+hv$ are exactly those of $v_3^2u_i^2u_j$, since $v^2$ has only
even powers. Dividing by $v_3$ before Cartier gives
$\sqrt{v_3b_j}\omega_i$.
For a degree-one double, $v$ has degree at most two, the norm
$v^2+hv+f=u_i^2u_j$ is monic of degree five, and
$g_i=(Y+v)/u_j$. The same calculation with $u_j$ monic linear
gives Cartier coefficient one.

Iteration gives $F^d(e_i)=A_i e_{qi}$, with $A_i$ as in (AC.14).
The hypothesis $\sigma_q(a)=qa$ makes coefficientwise conjugation
of the reduced representative $(u_i,v_i)$ equal to $(u_{qi},v_{qi})$.
Consequently $\sigma_q(\omega_i)=\omega_{qi}$ and
$\sigma_q(e_i)=e_{qi}$, with rigidifications unchanged.
The linearized coherent operator is $\sigma_q^{-1}F^d$,
whose scalar is $\sigma_q^{-1}(A_i)$.
This explains the $q$th-root correction; it cannot change whether
the scalar is one.

Finally, $\sigma_q(a)=qa$ is precisely invariance of the associated
geometric homomorphism to the constant group $\mathbf Z/n$.
The factor $q$ records Frobenius on $\mu_n$ in the Kummer
identification. The section of $P$ gives its unique pointed descent.
All multiplication functions were normalized at $P$, so no constant
character twist is omitted. Applying AC.17 to these one-dimensional
character spaces proves (AC.15). Selecting only primitive characters
gives the primitive-factor statement.
\end{proof}

